\documentclass[11pt,a4paper]{article}

% Tectonic = moteur XeTeX : UTF-8 natif, pas besoin de inputenc/fontenc T1
\usepackage[french]{babel}
\usepackage[a4paper, top=2.0cm, bottom=2.0cm, left=2.2cm, right=2.2cm]{geometry}
\usepackage{xcolor}
\usepackage{titlesec}
\usepackage{enumitem}
\usepackage{booktabs}
\usepackage{tabularx}
\usepackage{tcolorbox}
\usepackage[hidelinks,
    pdftitle={Plan de gestion des donnees - Rugulopteryx okamurae},
    pdfauthor={Jihan El-khattabi},
    pdfsubject={PGD / DMP OPIDoR - These de doctorat},
    pdfkeywords={PGD, DMP, OPIDoR, Rugulopteryx okamurae, phytochimie}
]{hyperref}

% Style des titres
\definecolor{accent}{HTML}{1F4E79}
\titleformat{\section}{\color{accent}\large\bfseries}{\thesection.}{0.5em}{}
\titlespacing*{\section}{0pt}{8pt}{4pt}
\setlist[itemize]{leftmargin=1.2em, itemsep=2pt, topsep=2pt}
\setlength{\parskip}{3pt}
\setlength{\parindent}{0pt}
\setlength{\emergencystretch}{2em}  % evite les overfull hbox sur les tokens \texttt

\begin{document}

% En-tete
\begin{center}
{\color{accent}\Large\bfseries Plan de gestion des données (PGD)}\\[2pt]
{\large\itshape Étude phytochimique de \textit{Rugulopteryx okamurae} dans la Méditerranée marocaine}\\[4pt]
\rule{0.95\linewidth}{0.4pt}\\[2pt]
\textbf{Modèle :} DMP OPIDoR (\url{https://dmp.opidor.fr}) -- format DMPRoadmap\\
\textbf{Doctorante :} Jihan El-khattabi \quad|\quad \textbf{ORCID :} \url{https://orcid.org/0009-0002-7946-4934}\\
\textbf{Laboratoire :} Biologie, Écologie et Santé (LBECS) \quad|\quad \textbf{Date :} 15/11/2026
\end{center}

\vspace{2pt}

\section*{1. Description des données collectées et produites}
La thèse génère plusieurs types de données issues de l'étude phytochimique de
\textit{Rugulopteryx okamurae}, prélevée sur la côte méditerranéenne marocaine
(détroit de Gibraltar et mer d'Alboran). Les données se répartissent en trois
ensembles : (i)~\textbf{données d'échantillonnage} -- fiches de terrain géoréférencées
(lat/lon GPS), date, profondeur, photographies des thalles ; (ii)~\textbf{données
analytiques} -- chromatogrammes HPLC-DAD, spectres de masse (HR-ESI-MS), spectres
RMN \textsuperscript{1}H et \textsuperscript{13}C, données d'activité biologique
(anti-inflammatoire, antioxydante) ; (iii)~\textbf{données dérivées} -- tables
d'identification des métabolites, structures chimiques, figures et cartes de
distribution. Le volume estimé est de l'ordre de 5 à 15~Go (essentiellement des
spectres RMN bruts en format FID), avec un volume réduit de tables traitées ($<$~50~Mo).

\section*{2. Normes, formats et métadonnées}
Les données brutes seront conservées dans leurs formats natifs : FID/ProcNMR
(Bruker), \texttt{.lcd}/\texttt{.lcdx} (HPLC Shimadzu), \texttt{.raw}/\texttt{.mzML}
pour la spectrométrie de masse, et \texttt{.csv} pour les tables. Les métadonnées
respecteront le standard \textbf{Darwin Core} pour les occurrences biologiques
(géoréférencement, identifiant taxonomique AlgaeBase) et le modèle \textbf{ISA-Tab}
pour les données métabolomiques. Chaque fichier sera documenté par un fichier
\texttt{README} compagnon décrivant l'auteur, la date, le protocole et les unités.
La nomenclature des fichiers suivra le schéma
\texttt{RO\_site\_date\_type\_n}, par exemple \texttt{RO\_Tetouan\_RMN\_001}.

\section*{3. Stockage, sécurité et sauvegarde}
\begin{itemize}
  \item \textbf{Stockage actif :} disque du laboratoire (serveur NAS, sauvegarde
        nocturne automatique) ; double local de travail sur disque dur externe chiffré.
  \item \textbf{Versionnement :} le code et les tables traitées sont versionnés dans
        le dépôt Git (\url{https://github.com/JihanElkhattabi/rugulopteryx-okamurae-these}),
        branche \texttt{main}, avec étiquettes de version (\texttt{v0.1}, \texttt{v1.0}).
  \item \textbf{Sécurité :} données sensibles (localisation précise des sites de
        prélèvement en cas de risque de biopiraterie) chiffrées (AES-256) ; accès
        restreint à la doctorante et au directeur de thèse.
  \item \textbf{Redondance :} règle 3-2-1 (3 copies, 2 supports, 1 hors site).
\end{itemize}

\section*{4. Partage, accès et dépôt long terme}
Après publication des résultats, les données seront déposées dans un entrepôt
\textbf{Nakala} (Huma-Num) pour les métadonnées et figures, et \textbf{Zenodo}
pour les données analytiques volumineuses, avec un DOI persistant. Les données
seront ouvertes sous licence \textbf{CC-BY-4.0}, dès l'acceptation des articles
associés, conformément à la politique de science ouverte du HCERES et au Plan
National pour la Science Ouverte. Un embargo maximum de 12~mois peut s'appliquer
aux données soumises à brevet ou à confidentialité industrielle (valorisation).
Les références bibliographiques sont partagées via la \textbf{bibliothèque Zotero}
de groupe publique.

\section*{5. Éthique, conformité légale et droits}
Les prélèvements d'algues sont réalisés dans le respect du \textbf{Code de la pêche
maritime} (loi n°~19-12) et des autorisations du Département de la Pêche Maritime ;
les zones protégées (parc national d'Al Hoceima) requièrent une autorisation du Haut
Commissariat aux Eaux et Forêts. Aucune donnée à caractère personnel n'est collectée,
hormis le consentement des partenaires de terrain. Les droits de propriété
intellectuelle appartiennent à l'université d'attache ; l'auteure conserve les droits
moraux. La citation de l'espèce suit Hwang, Lee \& Kim (2009) et Verlaque \textit{et al.} (2009).

\section*{6. Responsabilités et cycle de vie}
\begin{tabularx}{\linewidth}{lX}
\toprule
\textbf{Rôle} & \textbf{Responsabilité} \\
\midrule
Doctorante (J. El-khattabi) & Collecte, traitement, documentation, dépôt, charte IA \\
Directeur de thèse & Validation scientifique, autorisation de dépôt \\
Référent données du laboratoire & Conseil sur les standards, archivage Nakala/Zenodo \\
\bottomrule
\end{tabularx}\\[2pt]
Le cycle de vie prévu : collecte (2026-2027) $\rightarrow$ traitement (2027-2028)
$\rightarrow$ publication et dépôt (2028-2029) $\rightarrow$ archivage long terme
(min.~10~ans). Une révision annuelle du PGD est prévue pour intégrer les nouveaux
jeux de données et les outils du laboratoire.

\end{document}
